% !TeX encoding = UTF-8
% !TEX root = main.tex
% Catena di compilazione attuale: latexmk -> .latexmkrc -> lualatex -> pdf

\listfiles
\documentclass[a4paper, italian, 10pt, twoside, open=right]{scrbook}
\usepackage{z3f-preamble}% preambolo -> pacchetti custom -> pacchetti standard

%%%%%%%%%%%%%%%%%%%%
% NOTE DI UTILIZZO %
%%%%%%%%%%%%%%%%%%%%
% Nota 1: Aumentare Margini per rilegatura, bella l'efficienza nello spazio ma poco utile
% Nota 2: Trovare un'altra grafica per i bordi degli spazi di codice
% Nota 3: Cercare di rendere dinamico lo spazio di codice
% Nota 4: Prendere spunto dai template in giro per migliorare lo stile
% Nota 5: Riordinare i pacchetti in tematiche (ora che compilano)
% Nota 6: utilizzo di tectonic per la compilazione (rust)

%%%%%%%%%%%%
% SORGENTE %
%%%%%%%%%%%%

\begin{document}
\maketitle
\frontmatter
\hypersetup{linkcolor=black}
\setcounter{tocdepth}{2}
\tableofcontents
%\elencoteoremi
% \cleardoublepage
% \phantomsection
% \addcontentsline{toc}{chapter}{\listfigurename}
% \listoffigures

% \cleardoublepage
% \phantomsection
% \addcontentsline{toc}{chapter}{\listtablename}
% \listoftables

% \addcontentsline{toc}{chapter}{\lstlistlistingname}
% \lstlistoflistings

\mainmatter
\chapter{Introduzione}

Introduzione con del codice scritto a modino e pronto all'uso.

\begin{mintedcodebox}{c}{Codice di esempio di utilizzo della systemcall fork().}
#include <stdio.h>      /* for printf fprintf */
#include <sys/types.h>  /* for pid_t */
#include <unistd.h>     /* for fork() */
#include <stdlib.h>     /* for exit() */

int main(void) {

  pid_t pid;     /* variabile per memorizzare il PID del processo */

  pid = fork();  /* creazione processo figlio */

  /* ho due processi identici che continuano ad eseguire da qui */
  /* controllo il pid per distinguere il padre dal figlio */
  if(pid < 0){
    /* errore nell’esecuzione della fork */
    fprintf(stderr, "Generazione processo fallita...");
    exit(1);
  }else if (pid == 0){
    /* questo codice verra’ eseguito solo dal figlio */
    printf("Io sono il figlio con PID %d e PPID %d\n",
            getpid(), getppid());
  }else {
      /* questo codice lo esegue solo il processo padre */
      printf("Io sono il processo padre con PID %d\n", getpid());
      printf("Mio figlio ha PID %d\n", pid);
  }
  /* tutte le righe di codice che inserisco qua sarranno eseguite
    sia dal padre che dal figlio */
  printf("Processo con PID %d termina l’esecuzione\n", getpid());
  exit(0);
}
\end{mintedcodebox}


\end{document}
